\begin{document}
\lhead[ \leftmark ]{Informatik, Maturitätsarbeit}
\rhead[Informatik]{Cyril Wendl, \today, Winterthur}
\section*{Ideen für Maturitätsarbeiten im Fach Informatik}
Folgende Abschnitte stammen aus dem Vademecum der Maturitätsarbeit der Kantonsschule im Lee Winterthur (KLW). Sie sollen als Anregung für die Entwicklung von Themen und Projekten im Fach Informatik dienen. Es ist wichtig, dass die Schülerinnen und Schüler ihre eigenen Interessen und Ideen einbringen, um eine motivierende und erfolgreiche Maturitätsarbeit zu erstellen.
\begin{myExBox}
\subsection*{Eignung des Themas}

Das Interesse am Thema muss fast ein Jahr lang anhalten. Der erste spontane Einfall ist nicht
unbedingt der beste. Unreflektiert an alten Ideen zu kleben, kann einschränken und unflexibel
machen. Vor dem Entscheid sollten die Möglichkeiten und die Grenzen eines Themas durchdacht
werden.

Das Thema sollte nicht zu allgemein (zu offen) gehalten sein und auch nicht zu eng und
unergiebig. Ein konkret festgelegter Bereich (örtlich, zeitlich und thematisch fokussiert) und eine
Vorstellung, wie das Projekt methodisch angegangen werden könnte, ergeben zusammen die
Basis für eine gelingende Maturitätsarbeit.

\paragraph{Beispiele ungeeignete Themen (= zu allgemein und ungenau)}
\begin{itemize}
    \item Der Winterthurer Bahnhof
    \item Die Bachforelle
    \item Bild und Musik
\end{itemize}

\paragraph{Geeignete Themen (= konkret und präzise)}
\begin{itemize}
    \item Die Verkehrsplanung um den Winterthurer Bahnhof 2045+
    \item Untersuchung von Bachforellen-Standorten in einem Abschnitt der Töss
    \item Die Umsetzung von Miros Bild «Der Hund» in zeitgenössische Klaviermusik
\end{itemize}

\subsection*{Thema und Konzept entwickeln}

Eine Maturitätsarbeit hat gewisse Ansprüche zu erfüllen. Sie muss sowohl eigenständige,
kreative als auch analytische und reflektierende Anteile aufweisen und auf grundlegenden
Techniken, Methoden und Haltungen der Wissenschaft beruhen. Einfach nur ein Boot oder
Möbelstück bauen, einen Event organisieren oder eine Fotoserie erstellen, ist nicht ausreichend.

Das Thema ist in Absprache mit der Betreuungsperson so zu wählen, dass es inhaltlich sinnvoll
abgesteckt (allenfalls eingegrenzt oder erweitert) und sowohl methodisch als auch im gegebenen
Zeitrahmen geeignet bearbeitet werden kann.

Fächerübergreifende oder fachfremde Themen und Projekte sind grundsätzlich möglich, müssen
aber die Qualitätsansprüche einer Maturitätsarbeit und fachliche Standards erfüllen. In diesem
Fall sind die am besten geeigneten Richtlinien anderer Fachschaften anzuwenden.

\end{myExBox}
\newpage
\subsection*{Beispiele für mögliche Maturitätsarbeitsthemen im Fach Informatik}

\paragraph{1) Datenanalyse \& KI im Alltag}
\begin{itemize}
    \item \textbf{Wie gut erkennt ein KI-Textklassifikator Hassrede in deutschsprachigen Social-Media-Kommentaren?}\\
    Vergleich von regelbasiertem Ansatz und Transformer-Modell; Fokus auf Fehlklassifikationen, Fairness und ethische Konsequenzen.
    \item[\faStar] \textbf{Vom Klavier-Akkord zur Musiknote: Klassifikation von Klavierakkorden mit KI-Methoden}\\
    Aufbau eines Datensatzes von Klavierakkorden, Training verschiedener Klassifikatoren, Evaluierung der Genauigkeit und Diskussion von Anwendungsfällen.
\end{itemize}

\paragraph{2) Cybersicherheit \& Kryptologie}
\begin{itemize}
    \item[\faStar] \textbf{Phishing-Erkennung: Wirksamkeit unterschiedlicher Schulungsformen bei Jugendlichen}\\
    Konkretes Unterrichtssetting mit Vor-/Nachtest; statistische Auswertung der Erkennungsrate und Reflexion über nachhaltige Prävention.
    \item \textbf{Vergleich klassischer Verschlüsselung mit Post-Quanten-Verfahren in einem Messenger-Prototyp}\\
    Implementierung eines einfachen Prototyps, Messung von Laufzeit und Schlüssellängen, Diskussion von Praxistauglichkeit.
\end{itemize}

\paragraph{3) Softwareentwicklung \& Mensch-Computer-Interaktion}
\begin{itemize}
    \item[\faStar] \textbf{Entwicklung eines Passwortmanagers mit Python: Usability-Analyse und Sicherheitsbewertung}\\
    Vergleich von Verschlüsselungsalgorithmen, Implementierung eines Prototyps, Analyse der Sicherheitsaspekte.
    \item[\faStar] \textbf{LiveMacros: Entwicklung einer KI-App zur Visualisierung von Makronährstoffen}\\
    Prototypische Umsetzung, Nutzerbefragung zur Usability, Analyse von Stärken und Schwächen sowie Potenzialen für die Ernährungsbildung.
    \item[\faStar] \textbf{Von Augemented Reality zu Agents im Alltag: Wie verändert KI die unsere Interaktion mit der echten Welt?}\\
    Analyse von AR- und KI-Agenten mithilfe von Meta Glasses, Entwicklung einer einfachen AR-App, Diskussion von Chancen und Risiken für Bildung, Freizeit und Arbeit.
    \item[\faStar] \textbf{Winti wie Neu: Entwicklung einer Swift-App zur Erfassung von Nutzererfahrungen und Verbesserungsvorschlägen für die Winterthurer Innenstadt}\\
    Prototypische Umsetzung, Nutzerbefragung, Analyse der Daten und Ableitung von Handlungsempfehlungen für die Stadtplanung.
    \item[\faStar] \textbf{Programmierung eines 3D-Games mit Unreal Engine: Analyse von Entwicklungszeit, Performance und Nutzererlebnis}\\
    Entwicklung eines einfachen Spiels, Messung von Entwicklungszeit, FPS und Nutzerfeedback, Diskussion von Vor- und Nachteilen verschiedener Game-Plattformen.
\end{itemize}

\paragraph{4) Informatiksysteme \& Nachhaltigkeit}
\begin{itemize}
    \item \textbf{Energieverbrauch von Programmiersprachen bei identischen Algorithmen auf demselben Gerät}\\
    Vergleich z.\,B. Python, Java, C++; Messung von Laufzeit, CPU-Last und Energiebedarf; Diskussion methodischer Unsicherheiten.
    \item \textbf{Energiebedarf von KI-Systemen im schulischen Kontext}\\
    Analyse von Trainings- und Inferenzprozessen, Vergleich von lokalem vs. Cloud-basiertem Einsatz, Entwicklung von Richtlinien für nachhaltige Nutzung.
    \item \textbf{Optimierung von Bild- und Videodateien: Qualität über Quantität}\\
    Implementierung und Vergleich mehrerer Kompressionsverfahren mit objektiven Metriken (PSNR/SSIM) und subjektiver Qualitätsbeurteilung.
    \item \textbf{Rent a Human: Von der Sharing Economy zum Sharing von menschlicher Arbeitskraft}\\
    Analyse von Geschäftsmodellen, Nutzererfahrungen und gesellschaftlichen Auswirkungen, Entwicklung von Kriterien für faire und nachhaltige Plattformen im Zeitalter der KI-Agenten.
\end{itemize}

\paragraph{5) Robotik \& Physisches Computing}
\begin{itemize}
    \item[\faStar{}] \textbf{Autonomes Linienfolgen mit Hinderniserkennung: Vergleich von Roboter-Plattformern für den Informatikunterricht}\\
    Implementierung eines einfachen Prototyps auf Python, Messung von Genauigkeit und Reaktionszeit, Bewertung der Eignung für den Unterricht.
    \item \textbf{Indoor-Ortung im Klassenzimmer mit Bluetooth-Beacons}\\
    Aufbau eines Mini-Systems, Kalibrierung, Vergleich von Trilateration und Fingerprinting, Bewertung der Genauigkeit.
    \item \textbf{Intelligente Pflanzenstation mit Vorhersage des Giesszeitpunkts}\\
    Sensorik (Feuchte, Temperatur, Licht), Datenlogging über mehrere Wochen, Modellierung und Validierung im Realbetrieb.
\end{itemize}

\paragraph{6) Gesellschaft, Recht und Ethik der Informatik}
\begin{itemize}
    \item \textbf{Gesichtserkennung im öffentlichen Raum: technischer Leistungsvergleich und rechtliche Bewertung für die Schweiz}\\
    Kombination aus praktischen Tests (Fehlerraten unter variierenden Bedingungen) und Analyse rechtlicher Rahmenbedingungen.
    \item[\faStar] \textbf{Digitale Prüfungen mit KI-Hilfsmitteln: Welche Prüfungsformate bleiben fair und valide?}\\
    Entwicklung und Erprobung von Aufgabenformaten, Experteninterviews mit Lehrpersonen und begründete Handlungsempfehlungen.
\end{itemize}

\begin{myattention}[Hinweise zur Eingrenzung eines Themas]
Ein geeignetes MA-Thema wird präzise, wenn mindestens vier Elemente klar definiert sind:
\begin{itemize}
    \item \textbf{Kontext}: Wo findet die Untersuchung statt (z.\,B. konkrete Schule, Klasse, Datensatz, Strecke)?
    \item \textbf{Fragestellung}: Was genau soll beantwortet, verglichen oder optimiert werden?
    \item \textbf{Methode}: Wie werden Daten erhoben, ausgewertet und Ergebnisse validiert?
    \item \textbf{Ergebnisform}: Prototyp, Messresultate, Leitfaden, Kriterienkatalog oder Handlungsempfehlung.
\end{itemize}
\end{myattention}
\end{document}
